\documentclass[12pt]{article}

\newcommand{\en}{\selectlanguage{english}}
\newcommand{\gr}{\selectlanguage{greek}}

% Useful packages
\usepackage{caption}
\usepackage{amsmath}
\usepackage{amsfonts}
\usepackage{amssymb}
\usepackage{graphicx}
\usepackage{tabularray}
\usepackage{listings}
\usepackage{algorithm}
\usepackage{algpseudocode}
\usepackage[utf8]{inputenc}
\usepackage[T1]{fontenc}
\usepackage{csquotes}
\usepackage{enumitem}
\usepackage{mathtools} % Για το \coloneqq, \DeclarePairedDelimiter, κ.λπ.
\usepackage[a4paper,top=2cm,bottom=2cm,left=2cm,right=2cm,marginparwidth=1.75cm]{geometry}
\usepackage[english,main=greek]{babel}
\usepackage[dvipsnames]{xcolor}
\usepackage{pgfplots}
\usepackage{tikz}
\usetikzlibrary{trees}
\pgfplotsset{compat=newest}

\definecolor{col_bg}{HTML}{FFFFFF} 
\definecolor{col_txt}{HTML}{000000}
\pagecolor{col_bg}
\color{col_txt}


\begin{document}

\begin{table}[ht]
    \begin{tblr}{
      @{}X[l,valign=b]X[c,valign=b]X[r,valign=b]@{}
    }

        % First line, course info
        \hline
        \SetCell[c=2]{l}{Υπολογιστική Θεωρία Μηχανικής μάθησης} & & {2024-25} \\ 
        \hline
        {} & {} & {} \\

        % Title
        \SetCell[c=3]{c}{ \Large \textbf{Θεωρητικές ασκήσεις} } \\
        {} & {} & {} \\
        
        % Name Surname, RN
        \hline
        \SetCell[c=3]{c}{\textbf{Χολέβας Δημήτριος}} \\
        \hline
    \end{tblr}
\end{table}


\section*{1η Θεωρητική άσκηση (Εκτίμηση Κατάταξης)}

Αρχικά, αυτό που έχουμε να κάνουμε είναι να ορίσουμε $\text{$rank$}(x)$ ως τη θέση του $x$ στη φθίνουσα ταξινόμηση της λίστας. Στην ουσία θα ισχύει πως το μέγιστο στοιχείο έχει $rank$ $1$, το δεύτερο μέγιστο $rank$ $2$, κ.ο.κ. Ακολουθούμε τα παρακάτω βήματα:

\begin{enumerate}
  \item Λαμβάνουμε τυχαία ένα δείγμα μεγέθους $m$ από τη λίστα $x_1, \dots, x_n$, επιλέγοντας στοιχεία με επαναλήψεις (δηλαδή με αντικατάσταση).
  \item Ταξινομούμε τα δείγματα κατά φθίνουσα σειρά, ώστε τα μεγαλύτερα να έρχονται πρώτα.
  \item Εντοπίζουμε το στοιχείο που βρίσκεται στη θέση $\left\lfloor \frac{k m}{n} \right\rfloor$ στη ταξινομημένη υπολίστα και το επιστρέφουμε ως εκτίμηση.
\end{enumerate}
\subsection*{Ανάλυση Σφάλματος}

Αξιοποιούμε την ιδιότητα ότι η εμπειρική κατανομή που προκύπτει από ένα τυχαίο δείγμα προσεγγίζει την πραγματική κατανομή του πληθυσμού. Για να ποσοτικοποιήσουμε αυτή την απόκλιση, χρησιμοποιούμε την ανισότητα \textbf{$Dvoretzky–Kiefer–Wolfowitz (DKW)$}. Σύμφωνα με αυτήν, αν πάρουμε $m$ ανεξάρτητα δείγματα $X_1, \dots, X_m$ από κατανομή με συνάρτηση κατανομής $F$, τότε η εμπειρική κατανομή $F_m$ ικανοποιεί:

\[
\Pr\left[\sup_{x} |F_m(x) - F(x)| > \varepsilon\right] \leq 2 e^{-2 m \varepsilon^2}
\]

Αυτό μας εξασφαλίζει ότι η μέγιστη κατακόρυφη απόκλιση μεταξύ $F_m(x)$ και $F(x)$ είναι μικρή με πιθανότητα τουλάχιστον $1 - \delta$.

Η συνάρτηση $F(x)$ εκφράζει το ποσοστό των στοιχείων της αρχικής λίστας που είναι μικρότερα ή ίσα από $x$. Επομένως, η θέση $k$ αντιστοιχεί περίπου στο ποσοστό $k/n$ ως προς την πραγματική κατανομή. Αντίστοιχα, στο δείγμα μεγέθους $m$, η θέση $k \cdot m / n$ αποτελεί μια φυσική εκτίμηση αυτής της ποσοστιαίας μονάδας.

Στόχος μας είναι να εγγυηθούμε ότι το στοιχείο που επιλέγουμε από το δείγμα έχει κατάταξη που προσεγγίζει το $k$ στην αρχική λίστα. Στην ουσία θέλουμε να ισχύει το εξής:

\[
\left| \frac{\text{$rank$}(x)}{n} - \frac{k}{n} \right| \leq \varepsilon \cdot \frac{k}{n}
\Rightarrow
\left| \text{$rank$}(x) - k \right| \leq \varepsilon k
\]

Με αυτόν τον τρόπο σιγουρεύουμε ότι η κατάταξη του εκτιμώμενου στοιχείου βρίσκεται εντός του διαστήματος $[(1 - \varepsilon)k,\ (1 + \varepsilon)k]$ με υψηλή πιθανότητα.
\subsection*{Υπολογίζουμε το μέγεθος του δείγματος $m$}

Από την $DKW$, έχουμε:

\[
\Pr\left[\sup_{x} |F_m(x) - F(x)| \leq \varepsilon\right] \geq 1 - \delta
\Rightarrow
2 e^{-2 m \varepsilon^2} \leq \delta
\Rightarrow
m \geq \frac{1}{2 \varepsilon^2} \log \left(\frac{2}{\delta}\right)
\]
Με βάση το παραπάνω, παρατηρούμε πως αν επιλέξουμε:

\[
m = O\left(\frac{1}{\varepsilon^2} \log \frac{1}{\delta} \right)
\]

τότε με πιθανότητα τουλάχιστον $1 - \delta$, θα μας επιστραφεί ένα $x$ τέτοιο ώστε να ισχύει
\[
(1 - \varepsilon)k \leq \text{$rank$}(x) \leq (1 + \varepsilon)k
\]

\end{document}
